\begin{document}

\title{使用 LaTeX 编写文章}
\author{MetaBlog 示例}{}{example@example.com}

\begin{abstract}
MetaBlog 支持一个面向网页文章写作的 LaTeX 子集。这篇文章展示 Go 渲染器可以直接解析的常见结构，以及交给 LaTeXML 处理的表格和伪代码块。
\end{abstract}

\section{章节和文本样式}

可以使用 \verb|\section|、\verb|\subsection|、\verb|\subsubsection| 和 \verb|\subsubsubsection| 组织文章结构。

\subsection{行内样式}

行内文本可以使用 \textbf{粗体}、\textit{斜体}、\emph{强调}、\texttt{等宽字体}、\textsf{无衬线字体}、\textrm{衬线字体}、\textsc{小型大写}、\underline{下划线} 和 \textcolor{#256d85}{彩色文本}。

声明式样式也可以用在文本型参数中。例如，下面这个 tcb 标题使用了居中和较大字号：

\begin{tcb}[#e8f3f8][#256d85][#f7fbfd]{\Large\centering 带样式的 TCB 标题}
\texttt{tcb} 环境会渲染为可折叠文本框，包含标题栏、正文区域和左侧强调竖线。
\end{tcb}

\subsection{显式换行}

使用 \verb|\\| 可以显式换行，但不会创建新的段落。\\
这一行会出现在生成 HTML 的 \texttt{<br>} 后面。

\section{列表}

支持多层嵌套列表：

\begin{itemize}
  \item 站点内容
  \begin{itemize}
    \item 关于页面
    \item 文章页面
    \item 索引页面
  \end{itemize}
  \item 站点元数据
  \begin{enumerate}
    \item \texttt{config.toml}
    \item \texttt{articles.toml}
  \end{enumerate}
\end{itemize}

描述列表可以使用标签：

\begin{description}
  \item[Build] 生成静态 HTML 到输出目录。
  \item[Serve] 启动本地 HTTP 服务器用于预览。
  \item[Watch] 服务器运行时自动重建发生变化的文档。
\end{description}

\section{表格}

\texttt{table} 环境的外壳由 MetaBlog 解析，因此 caption、label 和交叉引用由站点生成器统一管理。真正复杂的 \texttt{tabular} 或 \texttt{tabularx} 主体会交给 LaTeXML 转换。

\begin{table}
\caption{\centering 示例站点常见文件}
\label{tab:example-files}
\begin{tabular}{lll}
\hline
文件 & 作用 & 是否必须 \\
\hline
\texttt{data/config.toml} & 站点标题、logo、icon 和分页配置 & 是 \\
\texttt{data/articles.toml} & 注册需要构建的文章 & 是 \\
\texttt{data/about_page/main.tex} & 关于页面源文件 & 否 \\
\texttt{data/custom_components/page_footing.tex} & 全站页尾组件 & 否 \\
\hline
\end{tabular}
\end{table}

表 \ref{tab:example-files} 展示了一个示例站点中最常见的配置文件。其代码如下：
\begin{minted}{LaTeX}
\begin{table}
\caption{\centering 示例站点常见文件}
\label{tab:example-files}
\begin{tabular}{lll}
\hline
文件 & 作用 & 是否必须 \\
\hline
\texttt{data/config.toml} & 站点标题、logo、icon 和分页配置 & 是 \\
\texttt{data/articles.toml} & 注册需要构建的文章 & 是 \\
\texttt{data/about_page/main.tex} & 关于页面源文件 & 否 \\
\texttt{data/custom_components/page_footing.tex} & 全站页尾组件 & 否 \\
\hline
\end{tabular}
\end{table}
\end{minted}


\section{伪代码}

\texttt{algorithm} 和 \texttt{algorithm*} 属于复杂块，会整体交给 LaTeXML 处理。适合在文章中放置论文风格的\texttt{algorithm2e}伪代码描述。

\begin{algorithm}[tbp]
    \small
    \caption{algorithm2e Example}
    \label{alg:alg_example}
    \SetKwInput{KwData}{Input}
    \SetKwInput{KwResult}{Output}
    \KwData{The Raw LaTeX algorithm block}
    \KwResult{Rendered pseudo code block}
    \SetKw{To}{to}
    \SetKw{Append}{append}
    \SetKw{Into}{into}
    \SetKw{Break}{break}
    \SetKw{Return}{return}
    \SetKwComment{EmptyLine}{ }{ }
    \SetNoFillComment
    \For{all block \(\in\) articles}{
    \If{block is algorithm(*) or /tabular(x)}{
        Render the block by \texttt{LaTeXML}\;
    }{
        Render the block by inside complier\;
    }
    }
    \Return All Blocks
\end{algorithm}

上面Algorithm~\ref{alg:alg_example}的伪代码描述了 MetaBlog 的文档级增量构建判断。

其代码如下：
\begin{minted}{LaTeX}
\begin{algorithm}[tbp]
    \small
    \caption{algorithm2e Example}
    \label{alg:alg_example}
    \SetKwInput{KwData}{Input}
    \SetKwInput{KwResult}{Output}
    \KwData{The Raw LaTeX algorithm block}
    \KwResult{Rendered pseudo code block}
    \SetKw{To}{to}
    \SetKw{Append}{append}
    \SetKw{Into}{into}
    \SetKw{Break}{break}
    \SetKw{Return}{return}
    \SetKwComment{EmptyLine}{ }{ }
    \SetNoFillComment
    \For{all block \(\in\) articles}{
    \If{block is algorithm(*) or /tabular(x)}{
        Render the block by \texttt{LaTeXML}\;
    }{
        Render the block by inside complier\;
    }
    }
    \Return All Blocks
\end{algorithm}
\end{minted}

\section{公式}

行内公式会交给 KaTeX 渲染，例如 $f(x)=x^2+\alpha$。

块级公式也会保留给 KaTeX：

\[
\forall x \in X,\quad f(x) \ge 0
\]

带编号的公式可以设置 label：

\begin{equation}
\label{eq:example-loss}
L(\theta)=\sum_{i=1}^{n}(y_i-\hat{y}_i)^2
\end{equation}

公式 \ref{eq:example-loss} 可以在普通文本中引用。

\section{链接}

可以使用 \verb|\url{https://example.com}| 显示完整 URL，也可以使用 \verb|\href{https://github.com}{GitHub}| 指定自定义链接文字。

\end{document}
